\section{Discussion}
\label{sec:discussion}

\subsection{What a non-separable objective changes in practice}

\ph{one or two paragraphs} The separable reading of degree selection --- spend
where the local error is largest --- is not a poor approximation to the right
answer; it optimizes a different functional, and the previous paper measured
that the two can order two policies oppositely. What this paper adds is that
the right functional is still tractable, because the coupling has a prefix--
suffix structure, and that the part a deployable rule cannot easily obtain is
not the weight but the sign. The sharpest way to say it is that all three
criteria are readings of one trajectory kernel: the force-level rule ignores
it, the sensitivity-weighted repair retains only its local diagonal
$\mathbf K(t,t)$, and only the third retains the cross-epoch coupling that
encodes cancellation. The weight was never the missing ingredient, because the
weight \emph{is} that diagonal.

\subsection{The resolution scale as a design constraint}

\ph{one paragraph} One length scale, $\pi r/N$, turns out to fix five otherwise
unrelated design choices:

\begin{enumerate}
  \item the sampling requirement on the accumulation grid, so that the signed
    integral is not aliased (Eq.~\eqref{eq:decorrelation});
  \item the resolution a surrogate for the omitted direction would have to
    reach, and therefore its cost (Section~\ref{sec:controller-direction});
  \item the probe refresh interval;
  \item the probe's share of the budget, which is the arc integral of its
    reciprocal (Eq.~\eqref{eq:probe-overhead-arc}); and
  \item the coherence horizon of the pilot arc, which is what forces the plan
    to be phase-indexed and the probe to be online.
\end{enumerate}

The rule of thumb worth carrying away is
\emph{nothing that varies faster than the truncation's own resolution scale can
be planned in advance} --- and the corollary, which is the expensive half:
whatever must therefore be measured in flight costs a synthesis every time the
scale turns over.

\subsection{Relation to goal-oriented adaptivity}

\ph{one paragraph} \citep{becker2001optimal}. The structural difference is that
the functional here is quadratic in an integral of the residual, so the
adjoint weight is not a scalar field and the resulting allocation is coupled
rather than pointwise.

\subsection{When would a practitioner use this?}

\ph{one paragraph} The controller costs a pilot arc, an offline planning pass,
and an online probe of \ph{pct} of the gravity budget. Two of those three
amortize and one does not, which decides the answer. The pilot arc and the plan
are paid once per trajectory family and can be reused across a Monte Carlo
ensemble or a covariance study; the probe is paid on every arc, every time.
A single arc therefore pays the full setup for one use and is the worst case;
\ph{quantify the break-even in arcs.} If the campaign lands in the regime where
the probe alone exceeds the margin, the honest recommendation is \Clite\ or the
constant degree, and Section~\ref{sec:campaign} says which.

\subsection{Limits}

\begin{itemize}
  \item \textbf{Model-relative.} Everything is measured against an adopted
    reference expansion, not against the true lunar field.
  \item \textbf{Linearization.} Both the benchmark and the planner are
    first-order. Amplitudes are not interpreted below \ph{km} perilune, where
    the previous campaign found the first-order amplitude wrong by up to a
    factor of two while the sign survived.
  \item \textbf{Arc length.} Calibrated and evaluated on seven-day arcs.
    \ph{state the sixty-day outcome.} A schedule planned for one arc length
    does not transfer to another, because the perilune descends and the dwell
    shifts.
  \item \textbf{Geometry.} \ph{state which populations the claim holds on.}
  \item \textbf{Certificate scope.} The gap of
    Section~\ref{sec:algorithm-certificate} bounds the distance to the best
    schedule of the linearized reference-arc problem, not of the propagated
    one.
  \item \textbf{Overhead is not small.} The online probe consumes \ph{pct} of
    the budget, an order of magnitude more than the switching cost and enough
    to decide the comparison on its own. It is charged in full
    (Eq.~\eqref{eq:matching}), but a reader should treat any margin narrower
    than it as contingent on the measured overhead rather than on the
    allocation idea.
  \item \textbf{Screening is not independent of the objective.} Candidates were
    screened by the same functional the benchmark optimizes
    (Section~\ref{sec:campaign}). The propagated sign check is what tests that
    transfer, and it is load-bearing rather than confirmatory.
  \item \textbf{No located optimum.} No continuous optimum in $\beta$, in $k$,
    in $\delta$, in $n_s$ or in $J$ is located anywhere in this paper; all are
    sampled on declared grids.
  \item \textbf{The benchmark and the controller are not claimed with equal
    force.} \Asign\ is derived from the trajectory-error functional itself and
    is certified against it, so the attainability result stands on the
    derivation. \Cplan\ is not claimed to be the best deployable schedule ---
    it is the first interpretable, auditable, low-complexity one. Three of its
    design choices are measured rather than argued and each has a declared
    alternative: the coordinate descent against a global solve on short
    instances (Section~\ref{sec:ablations}), the frozen offline plan against
    the horizon over which the coupling actually reaches
    (Table~\ref{tab:horizon}), and the band probe against a state-feedback
    policy fitted to \Asign\ itself. Where a measurement favours the
    alternative, it is reported as favouring the alternative.
\end{itemize}

\subsection{Future work}

\ph{short paragraph} Each item below is named with the measurement in this
paper that decides whether it is worth doing, so the list is a set of
conditional branches rather than a wish list.

\begin{itemize}
  \item \emph{Receding-horizon replanning}, if the horizon curve of
    Table~\ref{tab:horizon} has its knee well inside the arc or the coherence
    horizon of Section~\ref{sec:controller-split} is short.
  \item \emph{A reduced cancellation state.} The coupling enters only through
    $\mathbf A_i$, whose rank is at most three by construction and whose
    eigenvalue spread is reported in Section~\ref{sec:ablations}. If the trace
    concentrates in one or two directions, the accumulated state the controller
    must carry is that many scalars rather than a six-vector, and it can be
    updated in flight instead of being planned offline --- a cheaper controller
    than the one flown here, and the same reduction that makes a trellis
    formulation tractable.
  \item \emph{A switching penalty}, which turns the coordinate descent into a
    trellis problem and is only worth the state if that reduction holds.
  \item \emph{A vector-tail surrogate} in place of the probe, along the
    cost--retention curve of Section~\ref{sec:ablations}, if the measured probe
    overhead stays at the level reported here.
  \item \emph{A stronger solver} --- relax-and-round on the Frank--Wolfe
    iterate, or a global solve --- if the gap decomposition attributes the
    certified gap to the descent rather than to the relaxation.
  \item Allocation against an ensemble of initial states rather than one;
    transfer to other fields \citep{konopliv2016mars}; and propagating
    coefficient uncertainty into the degree-variance spectrum the amplitude
    estimate rests on.
\end{itemize}
